% part: intuitionistic-logic
% chapter: soundness-completeness

\documentclass[../../../include/open-logic-chapter]{subfiles}

\begin{document}

\olchapter{int}{sc}{বিশুদ্ধতা ও সম্পূর্ণতা}

\begin{editorial}
  এই অধ্যায়ে বচনমূলক স্বজ্ঞাবাদী যুক্তিবিদ্যার বিশুদ্ধতা ও
  সম্পূর্ণতার ফলগুলি একত্র করা হয়েছে। একটি ভূমিকা যোগ করা দরকার।
  সম্পূর্ণতার প্রমাণে প্রমাণযোগ্যতা-সংক্রান্ত কয়েকটি তথ্য ব্যবহৃত
  হয়েছে; সেগুলি কোথাও স্পষ্টভাবে উক্ত ও প্রমাণ করা উচিত।
\end{editorial}
% BN-SRC-761 TARGET-CORRECTION-BEGIN
\emph{প্রমাণের পূর্বপ্রস্তুতি:} পরের প্রমাণগুলিতে নিচের স্বজ্ঞাবাদী নিষ্পাদন-তথ্য ব্যবহৃত হয়। অনুমিতি বাড়ালে পুরোনো সসীম নিষ্পাদন বৈধ থাকে; প্রত্যেক নিষ্পাদন কেবল সসীমসংখ্যক অনুমিতি ব্যবহার করে। মিথ্যাতা নিষ্পাদিত হলে মিথ্যাতা থেকে সিদ্ধান্তের স্বতঃসিদ্ধ ও মোডাস পোনেন্সে যেকোনো সূত্র নিষ্পাদিত হয়। শর্তাধীন নিষ্পাদনের উপপাদ্য বলে: অতিরিক্ত অনুমিতি থেকে সিদ্ধান্ত পাওয়া তখনই সম্ভব, যখন আগের অনুমিতিগুলি থেকে ওই অনুমিতি-থেকে-সিদ্ধান্ত শর্তবচন পাওয়া যায়। স্বাভাবিক নিষ্পাদনে এটি শর্তবচন-প্রবর্তন; বিপরীত দিকে শর্তবচন-অপসারণ। স্বতঃসিদ্ধিক পদ্ধতিতে অতিরিক্ত অনুমিতিটিকে তার নিজের শর্তবচনে পাঠানো পরিচয়-নিষ্পাদন; অন্য অনুমিতি ও স্বতঃসিদ্ধের ক্ষেত্রে প্রথম শর্তবচন-স্বতঃসিদ্ধ, এবং মোডাস পোনেন্সের ধাপে দ্বিতীয় শর্তবচন-স্বতঃসিদ্ধ দিয়ে সসীম নিষ্পাদনের দৈর্ঘ্যে আরোহ চলে। পরিচয়-নিষ্পাদনে প্রথম স্বতঃসিদ্ধের দুই instance ও দ্বিতীয় স্বতঃসিদ্ধের এক instance যথেষ্ট। ফলে একই আগের অনুমিতি থেকে বিয়োজন এবং তার দুই অংশের প্রত্যেকটি আলাদাভাবে যোগ করে একই সিদ্ধান্ত পাওয়া গেলে, দুই শর্তবচন ও বিয়োজন-অপসারণের স্বতঃসিদ্ধে ওই সিদ্ধান্তই আগের অনুমিতি থেকে পাওয়া যায়। এই তথ্যগুলি ধ্রুপদি দ্বিনঞর্থকরণ-অপসারণের ওপর নির্ভর করে না।
% BN-SRC-761 TARGET-CORRECTION-END

\olimport{soundness-axd}
\olimport{soundness-nd}
\olimport{lindenbaum}
\olimport{canonical-model}
\olimport{truth-lemma}
\olimport{completeness-thm}
\olimport{decidability}

\OLEndChapterHook

\end{document}
